% TOP_PROBLEM: {"id": 3043, "title": "3-accessibility of Fibonacci numbers", "category_id": 2, "category": {"id": 2, "name": "combinatorics", "display_name": "Combinatorics", "description": "Counting problems, graph theory, discrete structures.", "slug": "combinatorics", "order_index": 2, "created_at": "2026-07-31T15:26:25.671Z"}, "status": "open", "problem_number": "OPG-1825", "set_id": 12}
% FIRST_POSTED: 2026-10-01
% ATTEMPT_STATUS: unresolved
% UnsolvedMath ID 3043; source catalogue code OPG-1825
% !TeX program = lualatex
% Research attempt by GPT-6 Astra Ultra; no claim of a new complete solution.
\documentclass[11pt,a4paper]{article}
\usepackage[margin=25mm]{geometry}
\usepackage{fontspec}
\setmainfont{lmroman10-regular.otf}[BoldFont=lmroman10-bold.otf,
 ItalicFont=lmroman10-italic.otf,BoldItalicFont=lmroman10-bolditalic.otf]
\usepackage{amsmath,amssymb,mathtools}
\usepackage{unicode-math}
\setmathfont{latinmodern-math.otf}
\AtBeginDocument{\let\setminus\smallsetminus}
\usepackage{xurl,hyperref}
\hypersetup{colorlinks=true,urlcolor=blue}
\setcounter{secnumdepth}{0}
\setlength{\parindent}{0pt}
\setlength{\parskip}{5pt}
\setlength{\emergencystretch}{3em}
\begin{document}
\section*{3-accessibility of Fibonacci numbers}\label{3043}
{\small UnsolvedMath ID 3043 · OPG-1825 · Combinatorics · OpenGarden}
\subsection{Definitions and mathematical statement}
Define $F_1=F_2=1$ and $F_{j+2}=F_{j+1}+F_j$ for $j\ge1$, and put
$\mathcal F=\{F_j:j\ge1\}\subseteq\mathbb N$. This is a set, so the two
initial occurrences of $1$ have no multiplicity significance. In particular,
$0$ is not an allowed gap.

Given a set $S\subseteq\mathbb N$, an $S$-difference sequence of length
$k$ is a strictly increasing sequence $x_1<\cdots<x_k$ of positive integers
such that $x_{i+1}-x_i\in S$ for every $1\le i<k$.
For a colouring $\chi:\mathbb N\to[r]$, the sequence is monochromatic if
$\chi(x_1)=\cdots=\chi(x_k)$. The set $S$ is called $r$-accessible when
for every such colouring and every integer $k\ge2$, a monochromatic
$S$-difference sequence of length $k$ exists.

Is $\mathcal F$ three-accessible? Written with its quantifiers, the target is
\[
 \forall\chi:\mathbb N\to\{1,2,3\}\quad\forall k\ge2\quad
 \exists x_1<\cdots<x_k:
 \ \chi(x_i)=\chi(x_1),\ \ x_{i+1}-x_i\in\mathcal F.
\]
The gap may vary from one step to the next, and Fibonacci values may be
reused as gaps. The terms $x_i$ themselves need not be Fibonacci numbers.
The length is arbitrary but finite; the statement does not require one
infinite monochromatic sequence. A negative answer must provide a
three-colouring and a fixed finite length which no monochromatic
Fibonacci-difference sequence reaches.
\subsection{Short English statement}
Does every three-colouring of the positive integers contain arbitrarily long monochromatic increasing sequences with Fibonacci gaps?
\subsection{Research attempt}
\paragraph{Scope and process.}
This is a GPT-6 Astra Ultra research attempt dated 1 October 2026, using
primary-source browsing and elementary mathematical checks. The canonical
definition above is retained. Results below are restricted results or
reductions, not a claim of a new solution to the entire catalogue question.
No formal proof assistant or external mathematical referee checked this text.


\paragraph{A material literature update.}
The historical source asks whether Fibonacci numbers are three-accessible.
On 27 September 2026, William J. Wesley posted a primary preprint [R1]
claiming a negative answer, using an aperiodic colouring coded by an
irrational rotation and a finite directed graph. This source was checked
on 1 October 2026. Thus the historical ``open'' metadata must not be
read as a fresh assertion that no resolution has been announced.
We do not claim authorship of that result or a complete independent
verification of its finite certificate. The independent work below
explains rigorously why periodic counterexample searches cannot succeed.

\paragraph{Every modulus divides some positive Fibonacci number.}
Include $F_0=0,F_1=1$ for this argument only; the allowed gap remains
positive. For any $q\ge1$, the recurrence on pairs modulo $q$ is
\[
 T(u,v)=(v,u+v),\qquad T^{-1}(u,v)=(v-u,u).
\]
It is a permutation of the finite set $(\mathbb Z/q\mathbb Z)^2$.
Therefore the orbit of $(0,1)$ returns to $(0,1)$ after some
positive number $h$ of steps, with $h\le q^2$. This is a return
to the starting pair, not just an eventual repeat elsewhere, because
$T$ is invertible. Its first coordinate gives $F_h\equiv0\pmod q$.
Also $F_h>0$, since $h\ge1$. In the case $q=1$ one may simply
take the gap $F_1=1$.

\paragraph{Periodic and eventually periodic colourings cannot obstruct.}
Suppose $\chi(n+q)=\chi(n)$ for all $n\ge N$, with any finite
number of colours. Choose a positive Fibonacci multiple $D=F_h$
of $q$ as above and any $a\ge N$. Then
\[
 a,\ a+D,\ a+2D,\ldots
\]
is an infinite increasing monochromatic sequence. Every consecutive
gap is the allowed positive Fibonacci number $D$, reused legally.
In particular this colouring contains a sequence of every finite
length requested in the definition. The result is stronger than the
three-colour statement within the eventual-periodicity subcase.

This also invalidates a common attempted construction: a finite
colouring of residues modulo $q$ might forbid all monochromatic gaps
in a tested finite Fibonacci list, but the first positive Fibonacci
multiple of $q$ creates a monochromatic arithmetic chain automatically.
Testing only Fibonacci numbers below a numerical cutoff can therefore
misclassify every genuinely periodic proposal.

\paragraph{A correct finite-search reformulation.}
Fix $k\ge2$. For each $N$, consider all three-colourings of $[N]$
containing no monochromatic Fibonacci-difference sequence of length
$k$. Restriction from $[N+1]$ to $[N]$ preserves this property.
If such colourings exist for every $N$, their restriction tree has
at most three children per vertex and nodes at every depth. An
infinite branch exists: recursively choose a child with descendants
at arbitrarily large depths, which exists because there are only
finitely many children. The resulting infinite colouring avoids
length $k$, since any forbidden finite sequence would lie in some
initial interval. Conversely an infinite avoiding colouring supplies
all finite restrictions. This is a rigorous compactness equivalence,
not an inference from one large finite computation.

\paragraph{Verification and relationship with the recent claim.}
The divisibility lemma uses only the recurrence, so duplicate initial
values $F_1=F_2=1$ cause no difficulty. The positive chain starts after
the eventual-periodic threshold and never returns below it. In the
compactness argument $k$ must be fixed while $N$ grows: allowing the
avoided length to grow with $N$ would not prove non-accessibility.
An irrational floor coding as in [R1] need not be periodic in $n$,
so the periodic-colouring theorem does not contradict its claimed
counterexample. The missing independent verification here is the
all-Fibonacci increment bound and the finite graph certificate for
that aperiodic construction, or an alternative explicit aperiodic
counterexample with a proved uniform length bound.

\paragraph{Final status of this attempt.}
\textbf{UNRESOLVED as an independent full verification.} The eventual-
periodicity theorem and compactness reduction are proved. A negative
resolution of the historical question is reported in [R1]; this file
neither certifies that complete proof nor represents the conjecture
as unchallenged by current literature.

\subsection{Sources}
\begin{itemize}
\item[S1] ULAM.AI and the UnsolvedMath Contributors, supplied problems.json, record 3043 (OPG-1825), OpenGarden collection. Catalogue identity and source transcription; CC BY 4.0.
\url{https://huggingface.co/datasets/ulamai/UnsolvedMath}.
\item[S2] Open Problem Garden, 3-accessibility of Fibonacci numbers. Original problem and discussion; the mathematical formulation below is expanded from the supplied source record.
\url{http://www.openproblemgarden.org/op/3_accessibility_of_fibonacci_numbers}.

\item[R1] William J. Wesley, \emph{The Fibonacci numbers are not
3-accessible}, arXiv:2609.33744v1, submitted 27 September 2026.
Primary preprint; abstract and proof structure checked 1 October 2026.
\url{https://arxiv.org/html/2609.33744v1}.

\end{itemize}
\end{document}
